\documentclass[12pt]{article}

\usepackage{geometry}
\usepackage{longtable}
\usepackage{hyperref}
\usepackage{array}

\geometry{margin=1in}

\title{
    \textbf{Snapshot 3 Objective: Checkpoint 2} \\
    JPL Lunar Detection Pipeline
}
\author{
    Software Engineering Tools Final Project \\
    Course: CS3338 \\
    Team Members: [Name 1], [Name 2], [Name 3], [Name 4]
}
\date{[Insert Date]}

\begin{document}

\maketitle
\newpage

\section{Objective}

The objective of Snapshot 3 is to add another major feature to the JPL Lunar Detection Pipeline. For this checkpoint, the new feature will be crater size filtering and detection history.

This snapshot also updates the project documents, Jira tasks, TestRail test cases, and workflow diagram.

\section{Reflection on Snapshot 2}

During Snapshot 2, the team added simulated boulder detection. The project was updated to include boulder mock data, new Jira tasks, and the first TestRail test cases.

Completed items from Snapshot 2 include:

\begin{itemize}
    \item Added boulder detection mock data.
    \item Updated SDD and SRS.
    \item Updated Design Specification and User Manual.
    \item Created Jira tasks and bugs for boulder detection.
    \item Created Snapshot 2 TestRail test cases.
    \item Updated the workflow diagram.
\end{itemize}

Snapshot 3 builds on this progress by adding more useful ways to view and organize detection results.

\section{New Feature: Crater Size Filtering}

The first new feature for Snapshot 3 is crater size filtering. This feature allows users to filter crater detections based on estimated crater size.

\subsection{Feature Description}

Crater size filtering helps users focus on certain types of craters. For example, a user may want to view only craters larger than 25 meters or only small craters below 10 meters.

For this project, crater size filtering will be simulated using mock data.

\subsection{Example Filter Options}

\begin{itemize}
    \item Show all craters.
    \item Show small craters less than 10 meters.
    \item Show medium craters between 10 and 50 meters.
    \item Show large craters greater than 50 meters.
\end{itemize}

\section{New Feature: Detection History}

The second new feature for Snapshot 3 is detection history. This feature allows users to view previous image processing runs.

\subsection{Feature Description}

The History Page will show previous uploaded images and their processing results. This makes the simulated system feel more complete and realistic.

\subsection{Example History Data}

\begin{longtable}{|p{1in}|p{1.7in}|p{1.5in}|p{1.2in}|}
\hline
\textbf{Image ID} & \textbf{File Name} & \textbf{Status} & \textbf{Detections} \\
\hline
IMG-001 & lunar\_south\_pole.jpg & Complete & 14 \\
\hline
IMG-002 & crater\_field.png & Complete & 22 \\
\hline
IMG-003 & mare\_region.jpg & Failed & 0 \\
\hline
\end{longtable}

\section{Goals for Snapshot 3}

The goals for Snapshot 3 are:

\begin{itemize}
    \item Add crater size filtering mock data.
    \item Add detection history mock data.
    \item Update the Results Page design.
    \item Update the History Page design.
    \item Update the SDD.
    \item Update the SRS.
    \item Update the Design Specification.
    \item Update the User Manual.
    \item Add new Jira tasks and bugs.
    \item Create Snapshot 3 TestRail test cases.
    \item Export or save the Snapshot 3 TestRail summary report.
    \item Update the workflow diagram.
\end{itemize}

\section{Planned Jira Tasks}

Example Jira tasks for Snapshot 3 include:

\begin{longtable}{|p{1.5in}|p{4.5in}|}
\hline
\textbf{Issue Type} & \textbf{Task} \\
\hline
Task & Add crater size filter options. \\
\hline
Task & Add mock detection history data. \\
\hline
Task & Update Results Page design. \\
\hline
Task & Update History Page design. \\
\hline
Task & Update SRS with new filtering requirements. \\
\hline
Task & Create TestRail test cases for crater filtering. \\
\hline
Bug & Crater filter does not update detection count. \\
\hline
Bug & History page shows failed scans as complete. \\
\hline
\end{longtable}

\section{TestRail Plan}

Snapshot 3 will include a second TestRail report.

Example test cases include:

\begin{longtable}{|p{1.2in}|p{2.4in}|p{2.4in}|}
\hline
\textbf{Test ID} & \textbf{Test Case} & \textbf{Expected Result} \\
\hline
TC-006 & Filter small craters & Only small crater detections are displayed. \\
\hline
TC-007 & Filter medium craters & Only medium crater detections are displayed. \\
\hline
TC-008 & Filter large craters & Only large crater detections are displayed. \\
\hline
TC-009 & View detection history & Previous scans appear on the History Page. \\
\hline
TC-010 & Open previous scan result & Selected previous scan displays its detection results. \\
\hline
\end{longtable}

\section{Documents to Update}

The following documents should be updated during Snapshot 3:

\begin{itemize}
    \item SDD.tex
    \item SRS.tex
    \item DesignSpec.tex
    \item UserManual.tex
    \item README.md
    \item Snapshot3\_Checkpoint2.tex
\end{itemize}

\section{Expected Deliverables}

By the end of Snapshot 3, the project should include:

\begin{itemize}
    \item Crater size filtering mock data.
    \item Detection history mock data.
    \item Updated SDD.
    \item Updated SRS.
    \item Updated Design Specification.
    \item Updated User Manual.
    \item Updated README.
    \item Jira tasks and bugs for Snapshot 3.
    \item Snapshot 3 TestRail report.
    \item Updated workflow diagram.
\end{itemize}

\section{Reflection Placeholder}

This section will be completed after Snapshot 3 is finished. The team will describe what was completed, what issues came up, and what should be finalized during Snapshot 4.

\end{document}